\documentclass[11pt]{article}
\usepackage[margin=1in]{geometry}
\usepackage{amsmath,amssymb,amsthm,hyperref}
\hypersetup{hidelinks}
\emergencystretch=2em
\newtheorem{theorem}{Theorem}
\newtheorem{lemma}{Lemma}
\title{Conjecture 00000002150:\\The kernel of the signless Laplacian}
\author{GitHub: gaochengzhecpu}
\date{}
\begin{document}
\maketitle

\begin{abstract}
For every finite undirected simple graph, the nullity of its signless
Laplacian is exactly the number of its bipartite connected components.
We construct a basis: each bipartite component contributes its alternating
sign vector, extended by zero outside that component. Nonbipartite
components contribute nothing. The full theorem and the basis construction
are verified in Lean with Mathlib.
\end{abstract}

\section{Statement and interpretation of the source}
Let $G=(V,E)$ be a finite undirected simple graph. Write $A$ for its
adjacency matrix, $D=\operatorname{diag}(\deg(v))$ for its degree matrix,
and $Q=D+A$ for its signless Laplacian. A connected component is bipartite
when its induced graph admits a proper coloring with two colors. Isolated
vertices are bipartite components. The zero eigenspace is $\ker Q$.

\begin{theorem}\label{thm:main}
Let $\mathcal B$ be the set of bipartite connected components of $G$.
Then
\[
\dim_{\mathbb R}\ker Q=|\mathcal B|.
\]
For each $C\in\mathcal B$, choose a bipartition
$C=C_+\sqcup C_-$ and define $s_C:V\to\mathbb R$ by
\[
s_C(v)=
\begin{cases}
1&v\in C_+,\\
-1&v\in C_-,\\
0&v\notin C.
\end{cases}
\]
The family $(s_C)_{C\in\mathcal B}$ is a basis of $\ker Q$.
\end{theorem}

The source's first assertion is the nullity identity above. Its second
phrase, ``the kernel of the count is an arbitrary bipartite base graph,''
does not define an additional graph construction. We give it the concrete
kernel interpretation supplied by the theorem: an arbitrary bipartite
connected component contributes exactly its alternating sign vector,
and arbitrary disjoint unions contribute their corresponding basis
vectors. We do not assign a separate meaning to an undefined ``base
graph'' operation. This reading, including the explicit basis, is stated
in the formal proposition. As usual for the adjacency-matrix assertion,
``general graphs'' means arbitrary finite undirected simple graphs;
infinite graphs, directed graphs, and unspecified loop conventions are
outside this statement.

\section{The proof}
\begin{lemma}\label{lem:ker}
For $x:V\to\mathbb R$,
\[
x^{\mathsf T}Qx=\sum_{\{u,v\}\in E}(x(u)+x(v))^2.
\]
Consequently, $Qx=0$ if and only if $x(u)=-x(v)$ for every edge
$\{u,v\}\in E$.
\end{lemma}
\begin{proof}
Each vertex term $\deg(u)x(u)^2$ contributes $x(u)^2$ once for each
incident edge. Each unordered edge contributes $2x(u)x(v)$ to
$x^{\mathsf T}Ax$. Combining the terms gives the identity. Equivalently,
the expression is one half of the sum over all ordered adjacent pairs.
If $Qx=0$, the nonnegative sum is zero, so each summand vanishes.
Conversely, if adjacent values are opposites, then
\[
(Qx)(u)=\deg(u)x(u)+\sum_{v\sim u}x(v)
=\deg(u)x(u)-\deg(u)x(u)=0
\]
for every vertex $u$.
\end{proof}

\begin{lemma}\label{lem:nonbip}
If $x\in\ker Q$ is nonzero at a vertex of a component $C$, then $C$
is bipartite. Hence $x$ vanishes identically on every nonbipartite
component.
\end{lemma}
\begin{proof}
Along an edge, Lemma~\ref{lem:ker} preserves the square of the value.
Along a walk it therefore preserves its absolute value. Starting from
a vertex where $x\ne0$, all values in $C$ are nonzero. Color a vertex
according to whether its value is positive or negative. Adjacent values
are opposites, so adjacent vertices get different colors. This is a
proper two-coloring of the entire induced component.
\end{proof}

\begin{proof}[Proof of Theorem~\ref{thm:main}]
For each bipartite component $C$, the vector $s_C$ has opposite values
at the ends of each edge of $C$. Any edge meeting $C$ lies entirely in
$C$, and an edge outside $C$ has both values zero. Thus
Lemma~\ref{lem:ker} gives $s_C\in\ker Q$.

The vectors are linearly independent because they have pairwise
disjoint supports and none is zero. More explicitly, for a chosen
vertex $r_C\in C$, evaluating $\sum_C a_Cs_C=0$ at $r_C$ gives
$a_Cs_C(r_C)=0$, hence $a_C=0$.

Now let $x\in\ker Q$ and set
\[
a_C=x(r_C)s_C(r_C)\qquad(C\in\mathcal B).
\]
For adjacent vertices in $C$, both $x$ and $s_C$ change sign. Their
product is therefore constant along walks in $C$. For every $v\in C$,
\[
x(v)s_C(v)=x(r_C)s_C(r_C)=a_C.
\]
Since $s_C(v)^2=1$ on $C$, this gives $x(v)=a_Cs_C(v)$.
On a nonbipartite component, $x$ vanishes by Lemma~\ref{lem:nonbip}.
Consequently
\[
x=\sum_{C\in\mathcal B}a_Cs_C.
\]
This proves spanning and the basis claim. Taking the dimension gives
the asserted nullity. If $V$ is empty, both the basis and its index set
are empty and the same argument gives dimension zero. Isolated vertices
contribute the usual coordinate vectors.
\end{proof}

\section{Formal verification and scope}
The file \texttt{lean/Main.lean} uses Mathlib's actual
\texttt{SimpleGraph}, adjacency and degree matrices,
$\mathbb R$-linear maps, connected components, proper colorings, bases,
and module dimension. Bipartiteness is proper two-colorability of the
induced graph on a connected component's support. It is not replaced
by a table, a numerical surrogate, or a stipulated nullity.

The formal proof establishes the quadratic identity, positive
semidefiniteness, the exact kernel edge relation, the coloring extracted
from a nonzero kernel value, the alternating sign vectors, their linear
independence, the explicit kernel expansion, and the actual kernel
basis. The main dimension theorem is
\begin{center}
\texttt{Signless.nullity\_eq\_bipartite\_components}.
\end{center}
The final theorem
\begin{center}
\texttt{Signless.conjecture}
\end{center}
proves a universally quantified proposition over all finite vertex
types and all simple graphs on them. That proposition includes both
the dimension identity and a basis whose values are the explicitly
defined component sign vectors. No substantive step in the stated real
finite-graph theorem is left to a paper-only argument.

For complex vectors, the same conclusion also follows by splitting a
kernel vector into real and imaginary parts, since $Q$ has real entries.
This optional complex-field extension is a paper observation; the
checked theorem is over $\mathbb R$.

The strict Lean check treats warnings as errors. The principal results
depend only on \texttt{propext}, \texttt{Classical.choice}, and
\texttt{Quot.sound}. No additional axioms or unproved placeholders are
used. Development and self-review were performed by a delegated AI
agent; the coordinating agent reviews the work separately. No external
independent review is claimed.

\section{Source and reproduction}
The exact frozen bilingual competition statement is included as
\texttt{SOURCE.md}; its SHA-256 and provenance are recorded in
\texttt{verification/SOURCE\_PROVENANCE.md}. The source is conjecture
\href{https://github.com/The-Last-Math-Competition/The-Last-Math-Competition/blob/main/conjectures/00000002150.md}{00000002150}
in the competition repository.

The project pins Lean 4.19.0 and Mathlib commit
\begin{center}
\texttt{c44e0c8ee63ca166450922a373c7409c5d26b00b}.
\end{center}
From the \texttt{lean} directory, obtain the pinned dependencies and
their Mathlib cache, then run \texttt{lake build} and
\begin{center}
\texttt{lake env lean -DwarningAsError=true Main.lean}.
\end{center}
The file ends with explicit axiom audits. Compile this document with
Tectonic or another LaTeX engine supporting the listed packages.
\end{document}
